% =====================================================================
% m2.tex -- Member 2 (Controller FSM & Tag Comparator) report sections
%   * Literature Review, subsection B (Cache Controller Design & FSM Architectures)
%   * Proposed Solution, subsection B (Controller FSM & Tag Comparator)
%
% For M5: % =====================================================================
% m2.tex -- Member 2 (Controller FSM & Tag Comparator) report sections
%   * Literature Review, subsection B (Cache Controller Design & FSM Architectures)
%   * Proposed Solution, subsection B (Controller FSM & Tag Comparator)
%
% For M5: % =====================================================================
% m2.tex -- Member 2 (Controller FSM & Tag Comparator) report sections
%   * Literature Review, subsection B (Cache Controller Design & FSM Architectures)
%   * Proposed Solution, subsection B (Controller FSM & Tag Comparator)
%
% For M5: % =====================================================================
% m2.tex -- Member 2 (Controller FSM & Tag Comparator) report sections
%   * Literature Review, subsection B (Cache Controller Design & FSM Architectures)
%   * Proposed Solution, subsection B (Controller FSM & Tag Comparator)
%
% For M5: \input{sections/m2} into the IEEEtran report.
% Figure path assumes the report lives in docs/.
% Citation keys used: smith1982, hennessy2017, hill1989, patterson2020
% (bibitems provided at the bottom of this file).
% =====================================================================


% ---------------------- BEGIN LITERATURE REVIEW B ---------------------
\subsection{Cache Controller Design and FSM Architectures}
\label{sec:litrev-m2}

The cache controller is the coordinating heart of the cache sub-system.
While the storage arrays (data, tag, and valid-bit RAMs) form the data
path, the controller orchestrates access sequencing, hit/miss detection,
CPU stall generation, and line refills~\cite{smith1982,hennessy2017}.
Classical cache controller design balances two competing objectives:
minimising the hit latency (which directly dictates the primary clock
cycle or pipeline stall penalty) and managing the multi-cycle latency of
main-memory transactions without corrupting architectural state.

In early microprocessor caches and classical single-issue architectures,
controllers were implemented as synchronous finite state machines
(FSMs)~\cite{smith1982}. Smith observed that cache control complexity
is strongly dictated by the write policy and allocation strategy.
A \emph{write-through} cache with \emph{no-write-allocate} policy
simplifies controller design dramatically: on a write hit, the updated
word is written into the cache line while simultaneously being forwarded
to memory; on a write miss, the cache remains untouched, and the write
is dispatched directly to main memory. Consequently, a write-through
controller requires no dirty bits and no block eviction (write-back)
states prior to allocating a new line. Conversely, a \emph{write-back}
controller requires tracking dirty status per block and executing a
two-phase memory handshake on misses that replace dirty lines (write-back
followed by refill)~\cite{hennessy2017}.

Hennessy and Patterson~\cite{hennessy2017,patterson2020} divide the
controller's operational cycle into distinct phases: address decoding, tag
comparison, response signaling, and refill allocation. A central design
decision is whether tag comparison occurs combinationally within the
request cycle or across discrete pipeline registers. In a direct-mapped
architecture, because each address maps to exactly one line, the index
can be presented to asynchronous-read RAMs immediately, allowing the tag
and valid bit to be compared combinationally in the same cycle. This
enables a single-cycle hit latency ($t_{\mathrm{hit}} = 1\text{ cycle}$),
which is foundational to achieving low Average Memory Access Time (AMAT):
\begin{equation}
\mathrm{AMAT} = t_{\mathrm{hit}} + (1 - h) \cdot t_{\mathrm{miss\_penalty}},
\label{eq:amat-m2}
\end{equation}
where $h$ is the hit rate. Any additional state transition introduced on
the critical hit path adds a mandatory cycle penalty to every load and
store, degrading overall processor throughput.

When a miss occurs, the controller must assert a stall signal to the CPU
pipeline, hold the request address and write data stable, and initiate a
refill sequence with the memory interface~\cite{hill1989}. The controller
FSM must therefore support a robust request/acknowledge handshake with
variable-latency DRAM controllers, transitioning through allocate, refill,
and update states before de-asserting the stall condition.
% ----------------------- END LITERATURE REVIEW B ----------------------


% ---------------------- BEGIN PROPOSED SOLUTION B ---------------------
\subsection{Controller Finite State Machine and Tag Comparator}
\label{sec:solution-m2}

Member 2 is responsible for the control path of the direct-mapped cache,
consisting of two modular Register-Transfer-Level (RTL) components:
the parameterised \texttt{tag\_comparator} and the \texttt{cache\_controller\_fsm}.

\subsubsection{Tag Comparator Architecture (\texttt{tag\_comparator.v})}
The tag comparator performs high-speed, purely combinational equality
checking between the tag stored in the cache line and the tag field of
the incoming CPU byte address. To prevent false hits on stale data after
power-on, reset, or line flushes, the comparator output is strictly gated
by the valid bit stored in \texttt{valid\_ram}:
\begin{equation}
\mathrm{hit} = \mathrm{rvalid} \;\land\; (\mathrm{rtag} == \mathrm{tag}).
\label{eq:hit-logic}
\end{equation}
In the default configuration, \texttt{tag} and \texttt{rtag} are 22 bits
wide ($32 - 6 - 4$). If $\mathrm{rvalid} = 0$, $\mathrm{hit}$ is driven to
$0$ immediately, guaranteeing that the first access to an uninitialised
line generates a compulsory (cold) miss even if the uninitialised tag bits
happen to match. The module is fully parameterised using the shared team
parameters (\texttt{ADDR\_WIDTH}, \texttt{CACHE\_SIZE}, \texttt{BLOCK\_SIZE}),
automatically re-deriving the comparator bit-width when sweeping cache
geometries during Phase~II.

\subsubsection{Controller FSM Architecture (\texttt{cache\_controller\_fsm.v})}
The cache controller FSM manages the access sequence and memory
handshakes according to a write-through, no-write-allocate policy.
Figure~\ref{fig:fsm_diagram} illustrates the state transition diagram,
and Table~\ref{tab:fsm_transitions} details the complete state-transition
and output control matrix.

\begin{figure}[t]
\centering
\includegraphics[width=\columnwidth]{figures/fsm_diagram.pdf}
\caption{State transition diagram of the Cache Controller FSM (Member 2).
Fast-path hits complete in a single cycle directly from \texttt{IDLE}. Read
misses traverse \texttt{COMPARE\_TAG} $\rightarrow$ \texttt{ALLOCATE} $\rightarrow$
\texttt{UPDATE} $\rightarrow$ \texttt{RESPOND}. Write misses skip \texttt{UPDATE}
under no-write-allocate policy.}
\label{fig:fsm_diagram}
\end{figure}

The FSM comprises five architecturally defined states:
\begin{enumerate}
    \item \textbf{IDLE (\texttt{3'b000})}: The resting state. When no access
          is requested ($\texttt{cpu\_req} = 0$), all write enables and
          memory requests remain de-asserted. When $\texttt{cpu\_req} = 1$
          and $\mathrm{hit} = 1$, the combinational fast-path asserts
          $\texttt{cpu\_ready} = 1$ in the same clock cycle, enabling
          single-cycle reads and single-cycle write-throughs without entering
          multi-cycle stall states.
    \item \textbf{COMPARE\_TAG (\texttt{3'b001})}: Entered on a miss from
          \texttt{IDLE}. The CPU address, write data, and read/write control
          signals are latched into internal holding registers to maintain
          signal stability throughout potential multi-cycle memory latency.
          On a miss, the FSM transitions directly to \texttt{ALLOCATE}.
    \item \textbf{ALLOCATE (\texttt{3'b010})}: The controller asserts
          $\texttt{mem\_req} = 1$ to initiate a main-memory transaction.
          For a read miss ($\texttt{cpu\_wr} = 0$), the block-aligned address
          $\{\mathrm{tag}, \mathrm{index}, 4'\mathrm{b}0000\}$ is driven to
          fetch the full 16-byte (128-bit) line. For a write miss, the target
          word address and write data are driven ($\texttt{mem\_wr} = 1$).
          The FSM remains in \texttt{ALLOCATE} until the memory interface
          asserts $\texttt{mem\_ready} = 1$.
    \item \textbf{UPDATE (\texttt{3'b011})}: Active solely for read misses.
          Upon memory burst arrival, the FSM asserts $\texttt{data\_we\_line} = 1$,
          $\texttt{tag\_we} = 1$, and $\texttt{valid\_we} = 1$ (with
          $\texttt{valid\_wdata} = 1$) for exactly one clock cycle, latching
          the refilled block into data RAM, updating the tag RAM, and marking
          the line valid. Write misses bypass this state entirely.
    \item \textbf{RESPOND (\texttt{3'b100})}: The controller asserts
          $\texttt{cpu\_ready} = 1$, signaling to the CPU that the requested
          read data is valid on $\texttt{cpu\_rdata}$ or that the write miss
          dispatch has terminated. The FSM then returns to \texttt{IDLE}.
\end{enumerate}

\begin{table*}[t]
\centering
\caption{Cache Controller FSM State Transition and Output Matrix}
\label{tab:fsm_transitions}
\begin{tabular}{llccccccccc}
\hline
\textbf{Current State} & \textbf{Next State Condition} & \textbf{Next State} & \texttt{cpu\_ready} & \texttt{cpu\_stall} & \texttt{tag\_we} & \texttt{valid\_we} & \texttt{data\_we\_word} & \texttt{data\_we\_line} & \texttt{mem\_req} & \texttt{mem\_wr} \\
\hline
\texttt{IDLE} & $\texttt{!cpu\_req}$ & \texttt{IDLE} & 0 & 0 & 0 & 0 & 0 & 0 & 0 & 0 \\
\texttt{IDLE} & $\texttt{cpu\_req} \land \mathrm{hit} \land \texttt{!cpu\_wr}$ (Read Hit) & \texttt{IDLE} & 1 & 0 & 0 & 0 & 0 & 0 & 0 & 0 \\
\texttt{IDLE} & $\texttt{cpu\_req} \land \mathrm{hit} \land \texttt{cpu\_wr}$ (Write Hit) & \texttt{IDLE} & 1 & 0 & 0 & 0 & 1 & 0 & 1 & 1 \\
\texttt{IDLE} & $\texttt{cpu\_req} \land \neg\mathrm{hit}$ (Miss) & \texttt{COMPARE\_TAG} & 0 & 1 & 0 & 0 & 0 & 0 & 0 & 0 \\
\hline
\texttt{COMPARE\_TAG} & $\mathrm{hit}$ & \texttt{RESPOND} & 0 & 1 & 0 & 0 & $\texttt{saved\_wr}$ & 0 & $\texttt{saved\_wr}$ & $\texttt{saved\_wr}$ \\
\texttt{COMPARE\_TAG} & $\neg\mathrm{hit}$ & \texttt{ALLOCATE} & 0 & 1 & 0 & 0 & 0 & 0 & 1 & $\texttt{saved\_wr}$ \\
\hline
\texttt{ALLOCATE} & $\texttt{!mem\_ready}$ & \texttt{ALLOCATE} & 0 & 1 & 0 & 0 & 0 & 0 & 1 & $\texttt{saved\_wr}$ \\
\texttt{ALLOCATE} & $\texttt{mem\_ready} \land \texttt{!saved\_wr}$ (Read Refill) & \texttt{UPDATE} & 0 & 1 & 0 & 0 & 0 & 0 & 0 & 0 \\
\texttt{ALLOCATE} & $\texttt{mem\_ready} \land \texttt{saved\_wr}$ (Write Miss) & \texttt{RESPOND} & 0 & 1 & 0 & 0 & 0 & 0 & 0 & 0 \\
\hline
\texttt{UPDATE} & 1 cycle (RAM write complete) & \texttt{RESPOND} & 0 & 1 & 1 & 1 & 0 & 1 & 0 & 0 \\
\hline
\texttt{RESPOND} & 1 cycle (access complete) & \texttt{IDLE} & 1 & 0 & 0 & 0 & 0 & 0 & 0 & 0 \\
\hline
\end{tabular}
\end{table*}

\subsubsection{Cycle-by-Cycle Operational Analysis}
The control timing for each operational case is outlined below:
\begin{itemize}
    \item \textbf{Read Hit (1 cycle):} $\texttt{cpu\_req} = 1$, $\texttt{cpu\_wr} = 0$.
          Address decoding and tag comparison occur combinationally. $\mathrm{hit} = 1$
          causes $\texttt{cpu\_ready}$ to assert high immediately. Data RAM
          outputs the selected 32-bit word on $\texttt{cpu\_rdata}$ in the same cycle.
          Total latency = $1\text{ cycle}$.
    \item \textbf{Read Miss (15 cycles total):} $\texttt{cpu\_req} = 1$, $\mathrm{hit} = 0$.
          Cycle~1: Decode/compare miss; $\texttt{cpu\_stall} = 1$.
          Cycle~2: Enter \texttt{ALLOCATE}; assert $\texttt{mem\_req} = 1$,
          $\texttt{mem\_addr} = \{\mathrm{tag}, \mathrm{index}, 4'\mathrm{b}0000\}$.
          Cycles~3--13: Main memory accesses first word (10 cycles) plus 3 burst words
          (1 cycle each) = 13 cycles.
          Cycle~14: Enter \texttt{UPDATE}; write 128-bit line to data RAM, tag to tag RAM,
          and set valid bit.
          Cycle~15: Enter \texttt{RESPOND}; assert $\texttt{cpu\_ready} = 1$.
          Total miss penalty = $14\text{ stall cycles}$ ($15\text{ cycles total}$).
    \item \textbf{Write Hit (1 cycle):} $\texttt{cpu\_req} = 1$, $\texttt{cpu\_wr} = 1$,
          $\mathrm{hit} = 1$. The FSM asserts $\texttt{data\_we\_word} = 1$, updating
          the targeted word in the cache line via read-modify-write, and simultaneously
          asserts $\texttt{mem\_req} = 1, \texttt{mem\_wr} = 1$ to forward the write
          to main memory. $\texttt{cpu\_ready} = 1$ is returned in the same cycle.
    \item \textbf{Write Miss (Variable memory latency):} $\texttt{cpu\_req} = 1$,
          $\texttt{cpu\_wr} = 1$, $\mathrm{hit} = 0$. Under the no-write-allocate
          policy, the cache line is not fetched into cache. The FSM dispatches
          the single word directly to memory ($\texttt{mem\_req} = 1, \texttt{mem\_wr} = 1$)
          and bypasses \texttt{UPDATE}, preserving the existing cache contents without
          pollution.
\end{itemize}
% ----------------------- END PROPOSED SOLUTION B ----------------------


% -------------- REFERENCES (merge into the shared list) ---------------
% \begin{thebibliography}{9}
% \bibitem{smith1982} A.~J. Smith, ``Cache memories,'' \emph{ACM Computing
%   Surveys}, vol.~14, no.~3, pp.~473--530, Sep. 1982.
% \bibitem{hill1989} M.~D. Hill and A.~J. Smith, ``Evaluating associativity
%   in CPU caches,'' \emph{IEEE Trans. Computers}, vol.~38, no.~12, pp.~1612--1630,
%   Dec. 1989.
% \bibitem{hennessy2017} J.~L. Hennessy and D.~A. Patterson, \emph{Computer
%   Architecture: A Quantitative Approach}, 6th~ed.\ Cambridge, MA, USA:
%   Morgan Kaufmann, 2017.
% \bibitem{patterson2020} D.~A. Patterson and J.~L. Hennessy, \emph{Computer
%   Organization and Design RISC-V Edition}, 2nd~ed.\ Morgan Kaufmann, 2020.
% \end{thebibliography}
 into the IEEEtran report.
% Figure path assumes the report lives in docs/.
% Citation keys used: smith1982, hennessy2017, hill1989, patterson2020
% (bibitems provided at the bottom of this file).
% =====================================================================


% ---------------------- BEGIN LITERATURE REVIEW B ---------------------
\subsection{Cache Controller Design and FSM Architectures}
\label{sec:litrev-m2}

The cache controller is the coordinating heart of the cache sub-system.
While the storage arrays (data, tag, and valid-bit RAMs) form the data
path, the controller orchestrates access sequencing, hit/miss detection,
CPU stall generation, and line refills~\cite{smith1982,hennessy2017}.
Classical cache controller design balances two competing objectives:
minimising the hit latency (which directly dictates the primary clock
cycle or pipeline stall penalty) and managing the multi-cycle latency of
main-memory transactions without corrupting architectural state.

In early microprocessor caches and classical single-issue architectures,
controllers were implemented as synchronous finite state machines
(FSMs)~\cite{smith1982}. Smith observed that cache control complexity
is strongly dictated by the write policy and allocation strategy.
A \emph{write-through} cache with \emph{no-write-allocate} policy
simplifies controller design dramatically: on a write hit, the updated
word is written into the cache line while simultaneously being forwarded
to memory; on a write miss, the cache remains untouched, and the write
is dispatched directly to main memory. Consequently, a write-through
controller requires no dirty bits and no block eviction (write-back)
states prior to allocating a new line. Conversely, a \emph{write-back}
controller requires tracking dirty status per block and executing a
two-phase memory handshake on misses that replace dirty lines (write-back
followed by refill)~\cite{hennessy2017}.

Hennessy and Patterson~\cite{hennessy2017,patterson2020} divide the
controller's operational cycle into distinct phases: address decoding, tag
comparison, response signaling, and refill allocation. A central design
decision is whether tag comparison occurs combinationally within the
request cycle or across discrete pipeline registers. In a direct-mapped
architecture, because each address maps to exactly one line, the index
can be presented to asynchronous-read RAMs immediately, allowing the tag
and valid bit to be compared combinationally in the same cycle. This
enables a single-cycle hit latency ($t_{\mathrm{hit}} = 1\text{ cycle}$),
which is foundational to achieving low Average Memory Access Time (AMAT):
\begin{equation}
\mathrm{AMAT} = t_{\mathrm{hit}} + (1 - h) \cdot t_{\mathrm{miss\_penalty}},
\label{eq:amat-m2}
\end{equation}
where $h$ is the hit rate. Any additional state transition introduced on
the critical hit path adds a mandatory cycle penalty to every load and
store, degrading overall processor throughput.

When a miss occurs, the controller must assert a stall signal to the CPU
pipeline, hold the request address and write data stable, and initiate a
refill sequence with the memory interface~\cite{hill1989}. The controller
FSM must therefore support a robust request/acknowledge handshake with
variable-latency DRAM controllers, transitioning through allocate, refill,
and update states before de-asserting the stall condition.
% ----------------------- END LITERATURE REVIEW B ----------------------


% ---------------------- BEGIN PROPOSED SOLUTION B ---------------------
\subsection{Controller Finite State Machine and Tag Comparator}
\label{sec:solution-m2}

Member 2 is responsible for the control path of the direct-mapped cache,
consisting of two modular Register-Transfer-Level (RTL) components:
the parameterised \texttt{tag\_comparator} and the \texttt{cache\_controller\_fsm}.

\subsubsection{Tag Comparator Architecture (\texttt{tag\_comparator.v})}
The tag comparator performs high-speed, purely combinational equality
checking between the tag stored in the cache line and the tag field of
the incoming CPU byte address. To prevent false hits on stale data after
power-on, reset, or line flushes, the comparator output is strictly gated
by the valid bit stored in \texttt{valid\_ram}:
\begin{equation}
\mathrm{hit} = \mathrm{rvalid} \;\land\; (\mathrm{rtag} == \mathrm{tag}).
\label{eq:hit-logic}
\end{equation}
In the default configuration, \texttt{tag} and \texttt{rtag} are 22 bits
wide ($32 - 6 - 4$). If $\mathrm{rvalid} = 0$, $\mathrm{hit}$ is driven to
$0$ immediately, guaranteeing that the first access to an uninitialised
line generates a compulsory (cold) miss even if the uninitialised tag bits
happen to match. The module is fully parameterised using the shared team
parameters (\texttt{ADDR\_WIDTH}, \texttt{CACHE\_SIZE}, \texttt{BLOCK\_SIZE}),
automatically re-deriving the comparator bit-width when sweeping cache
geometries during Phase~II.

\subsubsection{Controller FSM Architecture (\texttt{cache\_controller\_fsm.v})}
The cache controller FSM manages the access sequence and memory
handshakes according to a write-through, no-write-allocate policy.
Figure~\ref{fig:fsm_diagram} illustrates the state transition diagram,
and Table~\ref{tab:fsm_transitions} details the complete state-transition
and output control matrix.

\begin{figure}[t]
\centering
\includegraphics[width=\columnwidth]{figures/fsm_diagram.pdf}
\caption{State transition diagram of the Cache Controller FSM (Member 2).
Fast-path hits complete in a single cycle directly from \texttt{IDLE}. Read
misses traverse \texttt{COMPARE\_TAG} $\rightarrow$ \texttt{ALLOCATE} $\rightarrow$
\texttt{UPDATE} $\rightarrow$ \texttt{RESPOND}. Write misses skip \texttt{UPDATE}
under no-write-allocate policy.}
\label{fig:fsm_diagram}
\end{figure}

The FSM comprises five architecturally defined states:
\begin{enumerate}
    \item \textbf{IDLE (\texttt{3'b000})}: The resting state. When no access
          is requested ($\texttt{cpu\_req} = 0$), all write enables and
          memory requests remain de-asserted. When $\texttt{cpu\_req} = 1$
          and $\mathrm{hit} = 1$, the combinational fast-path asserts
          $\texttt{cpu\_ready} = 1$ in the same clock cycle, enabling
          single-cycle reads and single-cycle write-throughs without entering
          multi-cycle stall states.
    \item \textbf{COMPARE\_TAG (\texttt{3'b001})}: Entered on a miss from
          \texttt{IDLE}. The CPU address, write data, and read/write control
          signals are latched into internal holding registers to maintain
          signal stability throughout potential multi-cycle memory latency.
          On a miss, the FSM transitions directly to \texttt{ALLOCATE}.
    \item \textbf{ALLOCATE (\texttt{3'b010})}: The controller asserts
          $\texttt{mem\_req} = 1$ to initiate a main-memory transaction.
          For a read miss ($\texttt{cpu\_wr} = 0$), the block-aligned address
          $\{\mathrm{tag}, \mathrm{index}, 4'\mathrm{b}0000\}$ is driven to
          fetch the full 16-byte (128-bit) line. For a write miss, the target
          word address and write data are driven ($\texttt{mem\_wr} = 1$).
          The FSM remains in \texttt{ALLOCATE} until the memory interface
          asserts $\texttt{mem\_ready} = 1$.
    \item \textbf{UPDATE (\texttt{3'b011})}: Active solely for read misses.
          Upon memory burst arrival, the FSM asserts $\texttt{data\_we\_line} = 1$,
          $\texttt{tag\_we} = 1$, and $\texttt{valid\_we} = 1$ (with
          $\texttt{valid\_wdata} = 1$) for exactly one clock cycle, latching
          the refilled block into data RAM, updating the tag RAM, and marking
          the line valid. Write misses bypass this state entirely.
    \item \textbf{RESPOND (\texttt{3'b100})}: The controller asserts
          $\texttt{cpu\_ready} = 1$, signaling to the CPU that the requested
          read data is valid on $\texttt{cpu\_rdata}$ or that the write miss
          dispatch has terminated. The FSM then returns to \texttt{IDLE}.
\end{enumerate}

\begin{table*}[t]
\centering
\caption{Cache Controller FSM State Transition and Output Matrix}
\label{tab:fsm_transitions}
\begin{tabular}{llccccccccc}
\hline
\textbf{Current State} & \textbf{Next State Condition} & \textbf{Next State} & \texttt{cpu\_ready} & \texttt{cpu\_stall} & \texttt{tag\_we} & \texttt{valid\_we} & \texttt{data\_we\_word} & \texttt{data\_we\_line} & \texttt{mem\_req} & \texttt{mem\_wr} \\
\hline
\texttt{IDLE} & $\texttt{!cpu\_req}$ & \texttt{IDLE} & 0 & 0 & 0 & 0 & 0 & 0 & 0 & 0 \\
\texttt{IDLE} & $\texttt{cpu\_req} \land \mathrm{hit} \land \texttt{!cpu\_wr}$ (Read Hit) & \texttt{IDLE} & 1 & 0 & 0 & 0 & 0 & 0 & 0 & 0 \\
\texttt{IDLE} & $\texttt{cpu\_req} \land \mathrm{hit} \land \texttt{cpu\_wr}$ (Write Hit) & \texttt{IDLE} & 1 & 0 & 0 & 0 & 1 & 0 & 1 & 1 \\
\texttt{IDLE} & $\texttt{cpu\_req} \land \neg\mathrm{hit}$ (Miss) & \texttt{COMPARE\_TAG} & 0 & 1 & 0 & 0 & 0 & 0 & 0 & 0 \\
\hline
\texttt{COMPARE\_TAG} & $\mathrm{hit}$ & \texttt{RESPOND} & 0 & 1 & 0 & 0 & $\texttt{saved\_wr}$ & 0 & $\texttt{saved\_wr}$ & $\texttt{saved\_wr}$ \\
\texttt{COMPARE\_TAG} & $\neg\mathrm{hit}$ & \texttt{ALLOCATE} & 0 & 1 & 0 & 0 & 0 & 0 & 1 & $\texttt{saved\_wr}$ \\
\hline
\texttt{ALLOCATE} & $\texttt{!mem\_ready}$ & \texttt{ALLOCATE} & 0 & 1 & 0 & 0 & 0 & 0 & 1 & $\texttt{saved\_wr}$ \\
\texttt{ALLOCATE} & $\texttt{mem\_ready} \land \texttt{!saved\_wr}$ (Read Refill) & \texttt{UPDATE} & 0 & 1 & 0 & 0 & 0 & 0 & 0 & 0 \\
\texttt{ALLOCATE} & $\texttt{mem\_ready} \land \texttt{saved\_wr}$ (Write Miss) & \texttt{RESPOND} & 0 & 1 & 0 & 0 & 0 & 0 & 0 & 0 \\
\hline
\texttt{UPDATE} & 1 cycle (RAM write complete) & \texttt{RESPOND} & 0 & 1 & 1 & 1 & 0 & 1 & 0 & 0 \\
\hline
\texttt{RESPOND} & 1 cycle (access complete) & \texttt{IDLE} & 1 & 0 & 0 & 0 & 0 & 0 & 0 & 0 \\
\hline
\end{tabular}
\end{table*}

\subsubsection{Cycle-by-Cycle Operational Analysis}
The control timing for each operational case is outlined below:
\begin{itemize}
    \item \textbf{Read Hit (1 cycle):} $\texttt{cpu\_req} = 1$, $\texttt{cpu\_wr} = 0$.
          Address decoding and tag comparison occur combinationally. $\mathrm{hit} = 1$
          causes $\texttt{cpu\_ready}$ to assert high immediately. Data RAM
          outputs the selected 32-bit word on $\texttt{cpu\_rdata}$ in the same cycle.
          Total latency = $1\text{ cycle}$.
    \item \textbf{Read Miss (15 cycles total):} $\texttt{cpu\_req} = 1$, $\mathrm{hit} = 0$.
          Cycle~1: Decode/compare miss; $\texttt{cpu\_stall} = 1$.
          Cycle~2: Enter \texttt{ALLOCATE}; assert $\texttt{mem\_req} = 1$,
          $\texttt{mem\_addr} = \{\mathrm{tag}, \mathrm{index}, 4'\mathrm{b}0000\}$.
          Cycles~3--13: Main memory accesses first word (10 cycles) plus 3 burst words
          (1 cycle each) = 13 cycles.
          Cycle~14: Enter \texttt{UPDATE}; write 128-bit line to data RAM, tag to tag RAM,
          and set valid bit.
          Cycle~15: Enter \texttt{RESPOND}; assert $\texttt{cpu\_ready} = 1$.
          Total miss penalty = $14\text{ stall cycles}$ ($15\text{ cycles total}$).
    \item \textbf{Write Hit (1 cycle):} $\texttt{cpu\_req} = 1$, $\texttt{cpu\_wr} = 1$,
          $\mathrm{hit} = 1$. The FSM asserts $\texttt{data\_we\_word} = 1$, updating
          the targeted word in the cache line via read-modify-write, and simultaneously
          asserts $\texttt{mem\_req} = 1, \texttt{mem\_wr} = 1$ to forward the write
          to main memory. $\texttt{cpu\_ready} = 1$ is returned in the same cycle.
    \item \textbf{Write Miss (Variable memory latency):} $\texttt{cpu\_req} = 1$,
          $\texttt{cpu\_wr} = 1$, $\mathrm{hit} = 0$. Under the no-write-allocate
          policy, the cache line is not fetched into cache. The FSM dispatches
          the single word directly to memory ($\texttt{mem\_req} = 1, \texttt{mem\_wr} = 1$)
          and bypasses \texttt{UPDATE}, preserving the existing cache contents without
          pollution.
\end{itemize}
% ----------------------- END PROPOSED SOLUTION B ----------------------


% -------------- REFERENCES (merge into the shared list) ---------------
% \begin{thebibliography}{9}
% \bibitem{smith1982} A.~J. Smith, ``Cache memories,'' \emph{ACM Computing
%   Surveys}, vol.~14, no.~3, pp.~473--530, Sep. 1982.
% \bibitem{hill1989} M.~D. Hill and A.~J. Smith, ``Evaluating associativity
%   in CPU caches,'' \emph{IEEE Trans. Computers}, vol.~38, no.~12, pp.~1612--1630,
%   Dec. 1989.
% \bibitem{hennessy2017} J.~L. Hennessy and D.~A. Patterson, \emph{Computer
%   Architecture: A Quantitative Approach}, 6th~ed.\ Cambridge, MA, USA:
%   Morgan Kaufmann, 2017.
% \bibitem{patterson2020} D.~A. Patterson and J.~L. Hennessy, \emph{Computer
%   Organization and Design RISC-V Edition}, 2nd~ed.\ Morgan Kaufmann, 2020.
% \end{thebibliography}
 into the IEEEtran report.
% Figure path assumes the report lives in docs/.
% Citation keys used: smith1982, hennessy2017, hill1989, patterson2020
% (bibitems provided at the bottom of this file).
% =====================================================================


% ---------------------- BEGIN LITERATURE REVIEW B ---------------------
\subsection{Cache Controller Design and FSM Architectures}
\label{sec:litrev-m2}

The cache controller is the coordinating heart of the cache sub-system.
While the storage arrays (data, tag, and valid-bit RAMs) form the data
path, the controller orchestrates access sequencing, hit/miss detection,
CPU stall generation, and line refills~\cite{smith1982,hennessy2017}.
Classical cache controller design balances two competing objectives:
minimising the hit latency (which directly dictates the primary clock
cycle or pipeline stall penalty) and managing the multi-cycle latency of
main-memory transactions without corrupting architectural state.

In early microprocessor caches and classical single-issue architectures,
controllers were implemented as synchronous finite state machines
(FSMs)~\cite{smith1982}. Smith observed that cache control complexity
is strongly dictated by the write policy and allocation strategy.
A \emph{write-through} cache with \emph{no-write-allocate} policy
simplifies controller design dramatically: on a write hit, the updated
word is written into the cache line while simultaneously being forwarded
to memory; on a write miss, the cache remains untouched, and the write
is dispatched directly to main memory. Consequently, a write-through
controller requires no dirty bits and no block eviction (write-back)
states prior to allocating a new line. Conversely, a \emph{write-back}
controller requires tracking dirty status per block and executing a
two-phase memory handshake on misses that replace dirty lines (write-back
followed by refill)~\cite{hennessy2017}.

Hennessy and Patterson~\cite{hennessy2017,patterson2020} divide the
controller's operational cycle into distinct phases: address decoding, tag
comparison, response signaling, and refill allocation. A central design
decision is whether tag comparison occurs combinationally within the
request cycle or across discrete pipeline registers. In a direct-mapped
architecture, because each address maps to exactly one line, the index
can be presented to asynchronous-read RAMs immediately, allowing the tag
and valid bit to be compared combinationally in the same cycle. This
enables a single-cycle hit latency ($t_{\mathrm{hit}} = 1\text{ cycle}$),
which is foundational to achieving low Average Memory Access Time (AMAT):
\begin{equation}
\mathrm{AMAT} = t_{\mathrm{hit}} + (1 - h) \cdot t_{\mathrm{miss\_penalty}},
\label{eq:amat-m2}
\end{equation}
where $h$ is the hit rate. Any additional state transition introduced on
the critical hit path adds a mandatory cycle penalty to every load and
store, degrading overall processor throughput.

When a miss occurs, the controller must assert a stall signal to the CPU
pipeline, hold the request address and write data stable, and initiate a
refill sequence with the memory interface~\cite{hill1989}. The controller
FSM must therefore support a robust request/acknowledge handshake with
variable-latency DRAM controllers, transitioning through allocate, refill,
and update states before de-asserting the stall condition.
% ----------------------- END LITERATURE REVIEW B ----------------------


% ---------------------- BEGIN PROPOSED SOLUTION B ---------------------
\subsection{Controller Finite State Machine and Tag Comparator}
\label{sec:solution-m2}

Member 2 is responsible for the control path of the direct-mapped cache,
consisting of two modular Register-Transfer-Level (RTL) components:
the parameterised \texttt{tag\_comparator} and the \texttt{cache\_controller\_fsm}.

\subsubsection{Tag Comparator Architecture (\texttt{tag\_comparator.v})}
The tag comparator performs high-speed, purely combinational equality
checking between the tag stored in the cache line and the tag field of
the incoming CPU byte address. To prevent false hits on stale data after
power-on, reset, or line flushes, the comparator output is strictly gated
by the valid bit stored in \texttt{valid\_ram}:
\begin{equation}
\mathrm{hit} = \mathrm{rvalid} \;\land\; (\mathrm{rtag} == \mathrm{tag}).
\label{eq:hit-logic}
\end{equation}
In the default configuration, \texttt{tag} and \texttt{rtag} are 22 bits
wide ($32 - 6 - 4$). If $\mathrm{rvalid} = 0$, $\mathrm{hit}$ is driven to
$0$ immediately, guaranteeing that the first access to an uninitialised
line generates a compulsory (cold) miss even if the uninitialised tag bits
happen to match. The module is fully parameterised using the shared team
parameters (\texttt{ADDR\_WIDTH}, \texttt{CACHE\_SIZE}, \texttt{BLOCK\_SIZE}),
automatically re-deriving the comparator bit-width when sweeping cache
geometries during Phase~II.

\subsubsection{Controller FSM Architecture (\texttt{cache\_controller\_fsm.v})}
The cache controller FSM manages the access sequence and memory
handshakes according to a write-through, no-write-allocate policy.
Figure~\ref{fig:fsm_diagram} illustrates the state transition diagram,
and Table~\ref{tab:fsm_transitions} details the complete state-transition
and output control matrix.

\begin{figure}[t]
\centering
\includegraphics[width=\columnwidth]{figures/fsm_diagram.pdf}
\caption{State transition diagram of the Cache Controller FSM (Member 2).
Fast-path hits complete in a single cycle directly from \texttt{IDLE}. Read
misses traverse \texttt{COMPARE\_TAG} $\rightarrow$ \texttt{ALLOCATE} $\rightarrow$
\texttt{UPDATE} $\rightarrow$ \texttt{RESPOND}. Write misses skip \texttt{UPDATE}
under no-write-allocate policy.}
\label{fig:fsm_diagram}
\end{figure}

The FSM comprises five architecturally defined states:
\begin{enumerate}
    \item \textbf{IDLE (\texttt{3'b000})}: The resting state. When no access
          is requested ($\texttt{cpu\_req} = 0$), all write enables and
          memory requests remain de-asserted. When $\texttt{cpu\_req} = 1$
          and $\mathrm{hit} = 1$, the combinational fast-path asserts
          $\texttt{cpu\_ready} = 1$ in the same clock cycle, enabling
          single-cycle reads and single-cycle write-throughs without entering
          multi-cycle stall states.
    \item \textbf{COMPARE\_TAG (\texttt{3'b001})}: Entered on a miss from
          \texttt{IDLE}. The CPU address, write data, and read/write control
          signals are latched into internal holding registers to maintain
          signal stability throughout potential multi-cycle memory latency.
          On a miss, the FSM transitions directly to \texttt{ALLOCATE}.
    \item \textbf{ALLOCATE (\texttt{3'b010})}: The controller asserts
          $\texttt{mem\_req} = 1$ to initiate a main-memory transaction.
          For a read miss ($\texttt{cpu\_wr} = 0$), the block-aligned address
          $\{\mathrm{tag}, \mathrm{index}, 4'\mathrm{b}0000\}$ is driven to
          fetch the full 16-byte (128-bit) line. For a write miss, the target
          word address and write data are driven ($\texttt{mem\_wr} = 1$).
          The FSM remains in \texttt{ALLOCATE} until the memory interface
          asserts $\texttt{mem\_ready} = 1$.
    \item \textbf{UPDATE (\texttt{3'b011})}: Active solely for read misses.
          Upon memory burst arrival, the FSM asserts $\texttt{data\_we\_line} = 1$,
          $\texttt{tag\_we} = 1$, and $\texttt{valid\_we} = 1$ (with
          $\texttt{valid\_wdata} = 1$) for exactly one clock cycle, latching
          the refilled block into data RAM, updating the tag RAM, and marking
          the line valid. Write misses bypass this state entirely.
    \item \textbf{RESPOND (\texttt{3'b100})}: The controller asserts
          $\texttt{cpu\_ready} = 1$, signaling to the CPU that the requested
          read data is valid on $\texttt{cpu\_rdata}$ or that the write miss
          dispatch has terminated. The FSM then returns to \texttt{IDLE}.
\end{enumerate}

\begin{table*}[t]
\centering
\caption{Cache Controller FSM State Transition and Output Matrix}
\label{tab:fsm_transitions}
\begin{tabular}{llccccccccc}
\hline
\textbf{Current State} & \textbf{Next State Condition} & \textbf{Next State} & \texttt{cpu\_ready} & \texttt{cpu\_stall} & \texttt{tag\_we} & \texttt{valid\_we} & \texttt{data\_we\_word} & \texttt{data\_we\_line} & \texttt{mem\_req} & \texttt{mem\_wr} \\
\hline
\texttt{IDLE} & $\texttt{!cpu\_req}$ & \texttt{IDLE} & 0 & 0 & 0 & 0 & 0 & 0 & 0 & 0 \\
\texttt{IDLE} & $\texttt{cpu\_req} \land \mathrm{hit} \land \texttt{!cpu\_wr}$ (Read Hit) & \texttt{IDLE} & 1 & 0 & 0 & 0 & 0 & 0 & 0 & 0 \\
\texttt{IDLE} & $\texttt{cpu\_req} \land \mathrm{hit} \land \texttt{cpu\_wr}$ (Write Hit) & \texttt{IDLE} & 1 & 0 & 0 & 0 & 1 & 0 & 1 & 1 \\
\texttt{IDLE} & $\texttt{cpu\_req} \land \neg\mathrm{hit}$ (Miss) & \texttt{COMPARE\_TAG} & 0 & 1 & 0 & 0 & 0 & 0 & 0 & 0 \\
\hline
\texttt{COMPARE\_TAG} & $\mathrm{hit}$ & \texttt{RESPOND} & 0 & 1 & 0 & 0 & $\texttt{saved\_wr}$ & 0 & $\texttt{saved\_wr}$ & $\texttt{saved\_wr}$ \\
\texttt{COMPARE\_TAG} & $\neg\mathrm{hit}$ & \texttt{ALLOCATE} & 0 & 1 & 0 & 0 & 0 & 0 & 1 & $\texttt{saved\_wr}$ \\
\hline
\texttt{ALLOCATE} & $\texttt{!mem\_ready}$ & \texttt{ALLOCATE} & 0 & 1 & 0 & 0 & 0 & 0 & 1 & $\texttt{saved\_wr}$ \\
\texttt{ALLOCATE} & $\texttt{mem\_ready} \land \texttt{!saved\_wr}$ (Read Refill) & \texttt{UPDATE} & 0 & 1 & 0 & 0 & 0 & 0 & 0 & 0 \\
\texttt{ALLOCATE} & $\texttt{mem\_ready} \land \texttt{saved\_wr}$ (Write Miss) & \texttt{RESPOND} & 0 & 1 & 0 & 0 & 0 & 0 & 0 & 0 \\
\hline
\texttt{UPDATE} & 1 cycle (RAM write complete) & \texttt{RESPOND} & 0 & 1 & 1 & 1 & 0 & 1 & 0 & 0 \\
\hline
\texttt{RESPOND} & 1 cycle (access complete) & \texttt{IDLE} & 1 & 0 & 0 & 0 & 0 & 0 & 0 & 0 \\
\hline
\end{tabular}
\end{table*}

\subsubsection{Cycle-by-Cycle Operational Analysis}
The control timing for each operational case is outlined below:
\begin{itemize}
    \item \textbf{Read Hit (1 cycle):} $\texttt{cpu\_req} = 1$, $\texttt{cpu\_wr} = 0$.
          Address decoding and tag comparison occur combinationally. $\mathrm{hit} = 1$
          causes $\texttt{cpu\_ready}$ to assert high immediately. Data RAM
          outputs the selected 32-bit word on $\texttt{cpu\_rdata}$ in the same cycle.
          Total latency = $1\text{ cycle}$.
    \item \textbf{Read Miss (15 cycles total):} $\texttt{cpu\_req} = 1$, $\mathrm{hit} = 0$.
          Cycle~1: Decode/compare miss; $\texttt{cpu\_stall} = 1$.
          Cycle~2: Enter \texttt{ALLOCATE}; assert $\texttt{mem\_req} = 1$,
          $\texttt{mem\_addr} = \{\mathrm{tag}, \mathrm{index}, 4'\mathrm{b}0000\}$.
          Cycles~3--13: Main memory accesses first word (10 cycles) plus 3 burst words
          (1 cycle each) = 13 cycles.
          Cycle~14: Enter \texttt{UPDATE}; write 128-bit line to data RAM, tag to tag RAM,
          and set valid bit.
          Cycle~15: Enter \texttt{RESPOND}; assert $\texttt{cpu\_ready} = 1$.
          Total miss penalty = $14\text{ stall cycles}$ ($15\text{ cycles total}$).
    \item \textbf{Write Hit (1 cycle):} $\texttt{cpu\_req} = 1$, $\texttt{cpu\_wr} = 1$,
          $\mathrm{hit} = 1$. The FSM asserts $\texttt{data\_we\_word} = 1$, updating
          the targeted word in the cache line via read-modify-write, and simultaneously
          asserts $\texttt{mem\_req} = 1, \texttt{mem\_wr} = 1$ to forward the write
          to main memory. $\texttt{cpu\_ready} = 1$ is returned in the same cycle.
    \item \textbf{Write Miss (Variable memory latency):} $\texttt{cpu\_req} = 1$,
          $\texttt{cpu\_wr} = 1$, $\mathrm{hit} = 0$. Under the no-write-allocate
          policy, the cache line is not fetched into cache. The FSM dispatches
          the single word directly to memory ($\texttt{mem\_req} = 1, \texttt{mem\_wr} = 1$)
          and bypasses \texttt{UPDATE}, preserving the existing cache contents without
          pollution.
\end{itemize}
% ----------------------- END PROPOSED SOLUTION B ----------------------


% -------------- REFERENCES (merge into the shared list) ---------------
% \begin{thebibliography}{9}
% \bibitem{smith1982} A.~J. Smith, ``Cache memories,'' \emph{ACM Computing
%   Surveys}, vol.~14, no.~3, pp.~473--530, Sep. 1982.
% \bibitem{hill1989} M.~D. Hill and A.~J. Smith, ``Evaluating associativity
%   in CPU caches,'' \emph{IEEE Trans. Computers}, vol.~38, no.~12, pp.~1612--1630,
%   Dec. 1989.
% \bibitem{hennessy2017} J.~L. Hennessy and D.~A. Patterson, \emph{Computer
%   Architecture: A Quantitative Approach}, 6th~ed.\ Cambridge, MA, USA:
%   Morgan Kaufmann, 2017.
% \bibitem{patterson2020} D.~A. Patterson and J.~L. Hennessy, \emph{Computer
%   Organization and Design RISC-V Edition}, 2nd~ed.\ Morgan Kaufmann, 2020.
% \end{thebibliography}
 into the IEEEtran report.
% Figure path assumes the report lives in docs/.
% Citation keys used: smith1982, hennessy2017, hill1989, patterson2020
% (bibitems provided at the bottom of this file).
% =====================================================================


% ---------------------- BEGIN LITERATURE REVIEW B ---------------------
\subsection{Cache Controller Design and FSM Architectures}
\label{sec:litrev-m2}

The cache controller is the coordinating heart of the cache sub-system.
While the storage arrays (data, tag, and valid-bit RAMs) form the data
path, the controller orchestrates access sequencing, hit/miss detection,
CPU stall generation, and line refills~\cite{smith1982,hennessy2017}.
Classical cache controller design balances two competing objectives:
minimising the hit latency (which directly dictates the primary clock
cycle or pipeline stall penalty) and managing the multi-cycle latency of
main-memory transactions without corrupting architectural state.

In early microprocessor caches and classical single-issue architectures,
controllers were implemented as synchronous finite state machines
(FSMs)~\cite{smith1982}. Smith observed that cache control complexity
is strongly dictated by the write policy and allocation strategy.
A \emph{write-through} cache with \emph{no-write-allocate} policy
simplifies controller design dramatically: on a write hit, the updated
word is written into the cache line while simultaneously being forwarded
to memory; on a write miss, the cache remains untouched, and the write
is dispatched directly to main memory. Consequently, a write-through
controller requires no dirty bits and no block eviction (write-back)
states prior to allocating a new line. Conversely, a \emph{write-back}
controller requires tracking dirty status per block and executing a
two-phase memory handshake on misses that replace dirty lines (write-back
followed by refill)~\cite{hennessy2017}.

Hennessy and Patterson~\cite{hennessy2017,patterson2020} divide the
controller's operational cycle into distinct phases: address decoding, tag
comparison, response signaling, and refill allocation. A central design
decision is whether tag comparison occurs combinationally within the
request cycle or across discrete pipeline registers. In a direct-mapped
architecture, because each address maps to exactly one line, the index
can be presented to asynchronous-read RAMs immediately, allowing the tag
and valid bit to be compared combinationally in the same cycle. This
enables a single-cycle hit latency ($t_{\mathrm{hit}} = 1\text{ cycle}$),
which is foundational to achieving low Average Memory Access Time (AMAT):
\begin{equation}
\mathrm{AMAT} = t_{\mathrm{hit}} + (1 - h) \cdot t_{\mathrm{miss\_penalty}},
\label{eq:amat-m2}
\end{equation}
where $h$ is the hit rate. Any additional state transition introduced on
the critical hit path adds a mandatory cycle penalty to every load and
store, degrading overall processor throughput.

When a miss occurs, the controller must assert a stall signal to the CPU
pipeline, hold the request address and write data stable, and initiate a
refill sequence with the memory interface~\cite{hill1989}. The controller
FSM must therefore support a robust request/acknowledge handshake with
variable-latency DRAM controllers, transitioning through allocate, refill,
and update states before de-asserting the stall condition.
% ----------------------- END LITERATURE REVIEW B ----------------------


% ---------------------- BEGIN PROPOSED SOLUTION B ---------------------
\subsection{Controller Finite State Machine and Tag Comparator}
\label{sec:solution-m2}

Member 2 is responsible for the control path of the direct-mapped cache,
consisting of two modular Register-Transfer-Level (RTL) components:
the parameterised \texttt{tag\_comparator} and the \texttt{cache\_controller\_fsm}.

\subsubsection{Tag Comparator Architecture (\texttt{tag\_comparator.v})}
The tag comparator performs high-speed, purely combinational equality
checking between the tag stored in the cache line and the tag field of
the incoming CPU byte address. To prevent false hits on stale data after
power-on, reset, or line flushes, the comparator output is strictly gated
by the valid bit stored in \texttt{valid\_ram}:
\begin{equation}
\mathrm{hit} = \mathrm{rvalid} \;\land\; (\mathrm{rtag} == \mathrm{tag}).
\label{eq:hit-logic}
\end{equation}
In the default configuration, \texttt{tag} and \texttt{rtag} are 22 bits
wide ($32 - 6 - 4$). If $\mathrm{rvalid} = 0$, $\mathrm{hit}$ is driven to
$0$ immediately, guaranteeing that the first access to an uninitialised
line generates a compulsory (cold) miss even if the uninitialised tag bits
happen to match. The module is fully parameterised using the shared team
parameters (\texttt{ADDR\_WIDTH}, \texttt{CACHE\_SIZE}, \texttt{BLOCK\_SIZE}),
automatically re-deriving the comparator bit-width when sweeping cache
geometries during Phase~II.

\subsubsection{Controller FSM Architecture (\texttt{cache\_controller\_fsm.v})}
The cache controller FSM manages the access sequence and memory
handshakes according to a write-through, no-write-allocate policy.
Figure~\ref{fig:fsm_diagram} illustrates the state transition diagram,
and Table~\ref{tab:fsm_transitions} details the complete state-transition
and output control matrix.

\begin{figure}[t]
\centering
\includegraphics[width=\columnwidth]{figures/fsm_diagram.pdf}
\caption{State transition diagram of the Cache Controller FSM (Member 2).
Fast-path hits complete in a single cycle directly from \texttt{IDLE}. Read
misses traverse \texttt{COMPARE\_TAG} $\rightarrow$ \texttt{ALLOCATE} $\rightarrow$
\texttt{UPDATE} $\rightarrow$ \texttt{RESPOND}. Write misses skip \texttt{UPDATE}
under no-write-allocate policy.}
\label{fig:fsm_diagram}
\end{figure}

The FSM comprises five architecturally defined states:
\begin{enumerate}
    \item \textbf{IDLE (\texttt{3'b000})}: The resting state. When no access
          is requested ($\texttt{cpu\_req} = 0$), all write enables and
          memory requests remain de-asserted. When $\texttt{cpu\_req} = 1$
          and $\mathrm{hit} = 1$, the combinational fast-path asserts
          $\texttt{cpu\_ready} = 1$ in the same clock cycle, enabling
          single-cycle reads and single-cycle write-throughs without entering
          multi-cycle stall states.
    \item \textbf{COMPARE\_TAG (\texttt{3'b001})}: Entered on a miss from
          \texttt{IDLE}. The CPU address, write data, and read/write control
          signals are latched into internal holding registers to maintain
          signal stability throughout potential multi-cycle memory latency.
          On a miss, the FSM transitions directly to \texttt{ALLOCATE}.
    \item \textbf{ALLOCATE (\texttt{3'b010})}: The controller asserts
          $\texttt{mem\_req} = 1$ to initiate a main-memory transaction.
          For a read miss ($\texttt{cpu\_wr} = 0$), the block-aligned address
          $\{\mathrm{tag}, \mathrm{index}, 4'\mathrm{b}0000\}$ is driven to
          fetch the full 16-byte (128-bit) line. For a write miss, the target
          word address and write data are driven ($\texttt{mem\_wr} = 1$).
          The FSM remains in \texttt{ALLOCATE} until the memory interface
          asserts $\texttt{mem\_ready} = 1$.
    \item \textbf{UPDATE (\texttt{3'b011})}: Active solely for read misses.
          Upon memory burst arrival, the FSM asserts $\texttt{data\_we\_line} = 1$,
          $\texttt{tag\_we} = 1$, and $\texttt{valid\_we} = 1$ (with
          $\texttt{valid\_wdata} = 1$) for exactly one clock cycle, latching
          the refilled block into data RAM, updating the tag RAM, and marking
          the line valid. Write misses bypass this state entirely.
    \item \textbf{RESPOND (\texttt{3'b100})}: The controller asserts
          $\texttt{cpu\_ready} = 1$, signaling to the CPU that the requested
          read data is valid on $\texttt{cpu\_rdata}$ or that the write miss
          dispatch has terminated. The FSM then returns to \texttt{IDLE}.
\end{enumerate}

\begin{table*}[t]
\centering
\caption{Cache Controller FSM State Transition and Output Matrix}
\label{tab:fsm_transitions}
\begin{tabular}{llccccccccc}
\hline
\textbf{Current State} & \textbf{Next State Condition} & \textbf{Next State} & \texttt{cpu\_ready} & \texttt{cpu\_stall} & \texttt{tag\_we} & \texttt{valid\_we} & \texttt{data\_we\_word} & \texttt{data\_we\_line} & \texttt{mem\_req} & \texttt{mem\_wr} \\
\hline
\texttt{IDLE} & $\texttt{!cpu\_req}$ & \texttt{IDLE} & 0 & 0 & 0 & 0 & 0 & 0 & 0 & 0 \\
\texttt{IDLE} & $\texttt{cpu\_req} \land \mathrm{hit} \land \texttt{!cpu\_wr}$ (Read Hit) & \texttt{IDLE} & 1 & 0 & 0 & 0 & 0 & 0 & 0 & 0 \\
\texttt{IDLE} & $\texttt{cpu\_req} \land \mathrm{hit} \land \texttt{cpu\_wr}$ (Write Hit) & \texttt{IDLE} & 1 & 0 & 0 & 0 & 1 & 0 & 1 & 1 \\
\texttt{IDLE} & $\texttt{cpu\_req} \land \neg\mathrm{hit}$ (Miss) & \texttt{COMPARE\_TAG} & 0 & 1 & 0 & 0 & 0 & 0 & 0 & 0 \\
\hline
\texttt{COMPARE\_TAG} & $\mathrm{hit}$ & \texttt{RESPOND} & 0 & 1 & 0 & 0 & $\texttt{saved\_wr}$ & 0 & $\texttt{saved\_wr}$ & $\texttt{saved\_wr}$ \\
\texttt{COMPARE\_TAG} & $\neg\mathrm{hit}$ & \texttt{ALLOCATE} & 0 & 1 & 0 & 0 & 0 & 0 & 1 & $\texttt{saved\_wr}$ \\
\hline
\texttt{ALLOCATE} & $\texttt{!mem\_ready}$ & \texttt{ALLOCATE} & 0 & 1 & 0 & 0 & 0 & 0 & 1 & $\texttt{saved\_wr}$ \\
\texttt{ALLOCATE} & $\texttt{mem\_ready} \land \texttt{!saved\_wr}$ (Read Refill) & \texttt{UPDATE} & 0 & 1 & 0 & 0 & 0 & 0 & 0 & 0 \\
\texttt{ALLOCATE} & $\texttt{mem\_ready} \land \texttt{saved\_wr}$ (Write Miss) & \texttt{RESPOND} & 0 & 1 & 0 & 0 & 0 & 0 & 0 & 0 \\
\hline
\texttt{UPDATE} & 1 cycle (RAM write complete) & \texttt{RESPOND} & 0 & 1 & 1 & 1 & 0 & 1 & 0 & 0 \\
\hline
\texttt{RESPOND} & 1 cycle (access complete) & \texttt{IDLE} & 1 & 0 & 0 & 0 & 0 & 0 & 0 & 0 \\
\hline
\end{tabular}
\end{table*}

\subsubsection{Cycle-by-Cycle Operational Analysis}
The control timing for each operational case is outlined below:
\begin{itemize}
    \item \textbf{Read Hit (1 cycle):} $\texttt{cpu\_req} = 1$, $\texttt{cpu\_wr} = 0$.
          Address decoding and tag comparison occur combinationally. $\mathrm{hit} = 1$
          causes $\texttt{cpu\_ready}$ to assert high immediately. Data RAM
          outputs the selected 32-bit word on $\texttt{cpu\_rdata}$ in the same cycle.
          Total latency = $1\text{ cycle}$.
    \item \textbf{Read Miss (15 cycles total):} $\texttt{cpu\_req} = 1$, $\mathrm{hit} = 0$.
          Cycle~1: Decode/compare miss; $\texttt{cpu\_stall} = 1$.
          Cycle~2: Enter \texttt{ALLOCATE}; assert $\texttt{mem\_req} = 1$,
          $\texttt{mem\_addr} = \{\mathrm{tag}, \mathrm{index}, 4'\mathrm{b}0000\}$.
          Cycles~3--13: Main memory accesses first word (10 cycles) plus 3 burst words
          (1 cycle each) = 13 cycles.
          Cycle~14: Enter \texttt{UPDATE}; write 128-bit line to data RAM, tag to tag RAM,
          and set valid bit.
          Cycle~15: Enter \texttt{RESPOND}; assert $\texttt{cpu\_ready} = 1$.
          Total miss penalty = $14\text{ stall cycles}$ ($15\text{ cycles total}$).
    \item \textbf{Write Hit (1 cycle):} $\texttt{cpu\_req} = 1$, $\texttt{cpu\_wr} = 1$,
          $\mathrm{hit} = 1$. The FSM asserts $\texttt{data\_we\_word} = 1$, updating
          the targeted word in the cache line via read-modify-write, and simultaneously
          asserts $\texttt{mem\_req} = 1, \texttt{mem\_wr} = 1$ to forward the write
          to main memory. $\texttt{cpu\_ready} = 1$ is returned in the same cycle.
    \item \textbf{Write Miss (Variable memory latency):} $\texttt{cpu\_req} = 1$,
          $\texttt{cpu\_wr} = 1$, $\mathrm{hit} = 0$. Under the no-write-allocate
          policy, the cache line is not fetched into cache. The FSM dispatches
          the single word directly to memory ($\texttt{mem\_req} = 1, \texttt{mem\_wr} = 1$)
          and bypasses \texttt{UPDATE}, preserving the existing cache contents without
          pollution.
\end{itemize}
% ----------------------- END PROPOSED SOLUTION B ----------------------


% -------------- REFERENCES (merge into the shared list) ---------------
% \begin{thebibliography}{9}
% \bibitem{smith1982} A.~J. Smith, ``Cache memories,'' \emph{ACM Computing
%   Surveys}, vol.~14, no.~3, pp.~473--530, Sep. 1982.
% \bibitem{hill1989} M.~D. Hill and A.~J. Smith, ``Evaluating associativity
%   in CPU caches,'' \emph{IEEE Trans. Computers}, vol.~38, no.~12, pp.~1612--1630,
%   Dec. 1989.
% \bibitem{hennessy2017} J.~L. Hennessy and D.~A. Patterson, \emph{Computer
%   Architecture: A Quantitative Approach}, 6th~ed.\ Cambridge, MA, USA:
%   Morgan Kaufmann, 2017.
% \bibitem{patterson2020} D.~A. Patterson and J.~L. Hennessy, \emph{Computer
%   Organization and Design RISC-V Edition}, 2nd~ed.\ Morgan Kaufmann, 2020.
% \end{thebibliography}
